\documentclass[../main.tex]{subfiles}

\graphicspath{{\subfix{../images/}}}

\begin{document}

\clearpage
\thispagestyle{empty}
\vspace*{7cm}
\section{Esercitazioni di laboratorio}
\vspace*{1.5cm}

Procediamo ora presentando le seguenti esercitazioni di laboratorio svolte durante l'intera durata del corso. Le esercitazioni qui presentate (I - VIII) sono quelle proposte durante l'anno accademico 2025 - 2026 e potrebbero non essere aggiornata rispetto agli anni successivi. Infine invitiamo il lettore a verificare personalmente le soluzioni proposte in quanto potrebbero essere non completamente corrette e non essere presentate nella loro interezza. Tutte le soluzioni proposte sono riportate subito di seguito il testo dell'esercizio proposto.


\clearpage
\thispagestyle{empty}
\subsection{Esercitazione I}

In questa esercitazione di Controlli Automatici si affrontano e si mettono in pratica i concetti fondamentali per l’analisi e la simulazione della risposta dei sistemi dinamici di natura meccanica ed elettrica.

La prima parte del lavoro è dedicata alla costruzione del modello, partendo dalle equazioni differenziali che governano il sistema (ad esempio leggi di Newton per sistemi massa–molla–smorzatore o leggi di Kirchhoff per circuiti RLC). Tramite la trasformata di Laplace, tali equazioni vengono trasformate in relazioni algebriche, rendendo possibile il calcolo della funzione di trasferimento ($G(s)$), strumento centrale per rappresentare il sistema nel dominio di Laplace e studiarne poli, zeri, stabilità e dinamica.

Successivamente si procede con il calcolo analitico delle risposte nel tempo, determinando l’andamento di grandezze come posizione, velocità, corrente o tensione a fronte di ingressi tipici (gradino, impulso, rampa). L’inversione della trasformata di Laplace (anche mediante scomposizione in fratti semplici) permette di ottenere le espressioni esplicite della risposta temporale, distinguendo tra transitorio e regime permanente e valutando indici prestazionali (tempo di assestamento, sovraelongazione, errore a regime).

Infine, i risultati teorici vengono confrontati e verificati tramite simulazione (ad esempio con strumenti come MATLAB/Simulink o ambienti equivalenti), evidenziando la coerenza tra modello analitico e comportamento simulato. Questo approccio integrato consente di consolidare i metodi di modellazione e analisi dei sistemi lineari tempo-invarianti, fornendo una base solida per lo studio e la progettazione dei sistemi di controllo.


\vspace*{1cm}
\subsubsection{Simulazione della risposta di un sistema meccanico}
\label{sim1}

\subsubsection*{Soluzione esercizio \ref{sim1}}


\vspace*{1cm}
\subsubsection{Simulazione della risposta di un sistema elettrico}
\label{sim2}

\subsubsection*{Soluzione esercizio \ref{sim2}}

\vspace*{1cm}
\subsubsection{Calcolo delle funzioni di trasferimento}
\label{sim3}

\subsubsection*{Soluzione esercizio \ref{sim3}}

\vspace*{1cm}
\subsubsection{Calcolo analitico di risposte nel tempo di sistemi dinamici mediante antitrasformata di Laplace}
\label{sim4}

\subsubsection*{Soluzione esercizio \ref{sim4}}

\end{document}